\backhead{OBJECTIVE TYPE QUESTIONS}{khred}
\begin{exlist}
\item The reach of a 300 cm trigger is exceeded on\opt{the loam}{the sand}{the clay loam}{the silt loam}
\item 1 mm of water on 1 ha equals\opt{1 m$^3$}{10 m$^3$}{100 m$^3$}{1000 m$^3$}
\item The pore that is just full at pF 2 has a radius of about\opt{1.5 $\mu$m}{15 $\mu$m}{150 $\mu$m}{1.5 mm}
\item In the Hooghoudt equation the flow below the drains depends on\opt{$h^2$}{$d\,h$}{$L$}{$q^2$}
\item Christiansen's factor for many outlets is close to\opt{0.25}{0.35}{0.5}{1.0}
\end{exlist}
\objkey{1.\ (b),\quad 2.\ (b),\quad 3.\ (b),\quad 4.\ (b),\quad 5.\ (b)}
\dykband{EXERCISES IN SWAFL}
{\small Open SWAFL (desktop or online) with the benchmark project. The chapter of the \href{https://github.com/halilunaziru73-creator/Coupled-Water-Fluxes-Pore-Scale-Dynamics-Precision-Irrigation-Mulching-Feedback-Loops/raw/main/Software/v1.1.0/SWAFL_Desktop/docs/manual_v1.1.0/SWAFL_Manual_v1.1.0.pdf}{SWAFL manual} that explains each task is given in brackets; record your values in a table and compare them with what is expected (\emph{Check}).}\par\vspace{3pt}
\begin{exlist}
\item {\color{khblue}[Manual, Chapter 6]} For each of the ten soils set a 300 cm trigger and record whether the crop is irrigated; compare with the table of this chapter.\par\noindent{\small\itshape Check: sand and loamy sand are not irrigated.}
\item {\color{khblue}[Manual, Chapter 11]} Check the surface design of the loam against the worked solution of this chapter.\par\noindent{\small\itshape Check: the values agree.}
\item {\color{khblue}[Manual, Chapter 13A]} Check the drain spacing of the worked example in the drainage layer.\par\noindent{\small\itshape Check: the spacing agrees within rounding.}
\item {\color{khblue}[Manual, Chapter 12]} Check the lateral loss of the worked example in the drip layer.\par\noindent{\small\itshape Check: the loss agrees within rounding.}
\end{exlist}
